\section{Reading Xie24}


\subsection{Setup}

  Let $\kkk$ be an algebraically closed field of arbitrary characteristic.
  Let $X$ be a projective variety of dimension $d$ over $\kkk$, and $f: X \dashrightarrow X$ be a dominant rational map.
  All birational models of $X$ form a directed system, we denote by $\frakX$.
  We define the inductive limit of the numerical groups $\upN^1(X')$ by 
  \[ \upN^1(\frakX) \coloneqq \upN^1(\frakX,\pi^*) \coloneqq \varinjlim_{\pi:X'' \to X' \in \frakX} \left( \pi^* \upN^1(X') \to \upN^1(X'') \right). \]
  One can view $\upN^1(\frakX)$ as the group of numerical classes of $\bfb$-Cartier divisors with $\bbR$-coefficients on $X$.
  For every $D \in \upN^1(\frakX)$, there exists a birational model $X' \in \frakX$ and $D' \in \upN^1(X')$ such that $D_{X''} = \pi^* D'$ for every $\pi: X'' \to X'$ in $\frakX$.
  
  
  % picard groups $\Pic(X')$ for $X' \in \frakX$ by
  % \[ \Pic(\frakX) := \varinjlim_{\pi:X'' \to X' \in \frakX} \left( \pi^* \Pic(X') \to \Pic(X'') \right). \]

  \begin{theorem}[Siu's inequality]
    Let $X$ be a projective variety of dimension $d$, and let $L,M$ be nef $\bbR$-divisors on $X$ with $M$ big. 
    Then 
    \[ (M^d) L \leq d (M^{d-1} \cdot L) M. \]
    In particular, for every $\varepsilon \in (0, 1)$, we have
    \[ d (M^{d-1} \cdot L) M - \varepsilon (M^d) L \]
    is big.
  \end{theorem}

\subsection{Results}

  Fix a projective variety $X$ of dimension $d$ and a dominant rational map $f: X \dashrightarrow X$. 

  \begin{mainthm}\label{thm:recursive_inequality}
    For $i = 0, 1, \ldots, d$ and $\varepsilon \in (0, 1)$, there exists $m_\varepsilon > 0$ such that for all $m \geq m_\varepsilon$,
    \[  L_{2m} - \varepsilon^m \mu_i^m L_m + \mu_i^m \mu_{i+1}^m L \]
    is big.
  \end{mainthm}

  \begin{mainthm}\label{mainthm:semi-continuity}
    Let $S$ be an integral noetherian scheme, and $\pi: \calX \to S$ be a flat projective morphism. 
    Let $f: \calX \dashrightarrow \calX$ be a dominant rational map over $S$.
    Then for every $i$, the function $s \mapsto \lambda_i(f_s)$ is lower semi-continuous.
  \end{mainthm}

  \begin{mainthm}
    If $f$ is cohomologically hyperbolic, then the set of periodic points of $f$ is Zariski dense in $X$.
  \end{mainthm}


\subsection{Inequalities}


  The main goal is to find some lower bound $\beta_m \leq \lambda_i$ of $\lambda_i$ such that $\beta_m \to \lambda_i$ as $m \to \infty$ and $\beta_m$ depends only on the intersection numbers of $L_{2m}, L_m, L$.


  Fix number $\alpha_i$, $\varepsilon, \delta \in (0, 1)$, and $m \in \bbZ_{>0}$. \Yang{here $m$ is to iterate $f$, hence if we allow iteration, we can drop the $m$ in the statement of \cref{thm:recursive_inequality}.}
  We have the following conditions:
  \begin{enumerate}
    \item condition $I_i$:
      \[ \frac{\bigl((f^*)^2L+\alpha_{j+1}\alpha_{i}\delta L\bigr)^{d-i+j+1} \cdot \bigl(f^*L\bigr)^{i-j-1}}{\alpha_{j+1}\bigl((f^*)^2L+\alpha_{j+1}\alpha_{i}\delta L\bigr)^{d-i+j} \cdot \bigl(f^*L\bigr)^{i-j}} > (d-i-j+1) \varepsilon; \]
      or the following is enough:
      \[ L_2 + \alpha_{j+1}\alpha_{i}\delta L >_{\text{big}} \varepsilon \alpha_{j+1} L_1. \]
    \item condition $J_i$:
  \end{enumerate}


  The logic chain is as follows:
  \[ I,J,K \implies \lambda_i \geq \beta_i \]
  and 
  \[ \mu_i > \mu_{i+1} \implies I,J,K \text{ hold for some } \mu_i. \]

  $\approx$ $\asymp$


  \begin{proposition}
    \Yang{}
    Let $m_i \geq 0$, $r_i > r_{i-1}$ with $i = 0, 1, \ldots, d$.
    Then there exists some constant $C > 0$ depending only on $(X,f,L)$ and $r_i$'s such that for every $m_i$'s, we have
    \[ C^{-1} \deg_d(f^{\sum m_i}) \prod_i \frac{\deg_{r_i}(f^{m_i})}{\deg_{r_{i-1}}(f^{m_i})} \leq (L_{M_1}^{r_1} \cdot L_{M_2}^{r_2 - r_1} \cdots L_{M_d}^{r_d - r_{d-1}}) \leq C \prod_i \deg_{r_i}(f^{m_i}). \]
  \end{proposition}

  Set 
  \[ L_m := (f^m)^* L, \quad M_m(t) := L_{2m} + t^m L. \]
  We will try to estimate the ratio
  \[ \gamma_{L,t}(m) = \frac{M_m(t)^{d}}{M_m(t)^{d-1} \cdot L_m}. \]






% \subsection{Semi-continuity}

%   Here we test \cref{mainthm:semi-continuity} by considering the family of 


% \subsection{Density of periodic points}




% \subsection{Other results}

%   \begin{theorem}
%     We have 
%     \[ \lambda_1(f) \geq \frac{\deg_{1,L} f^2}{4 \deg_{1,L} f}. \]
%   \end{theorem}

%   \begin{theorem}
%     \Yang{There is an algorithm to compute $\lambda_i$'s.}
%   \end{theorem}

%   \begin{theorem}
%     Let $X$ be a projective variety of dimension $d$, and $f: X \dashrightarrow X$ be a dominant rational map. 
%     Then for every $i = 0, 1, \ldots, d-1$, we have $\lambda_i(f) \leq \lambda_1(f)^i$.

%     Suppose that $\mu_2 < 1$ or $\mu_1 = \mu_2$.
%   \end{theorem}